\section*{Capítulo \ref{ch:b1:arith} --- Aritmética dos Inteiros}

\begin{solution}{exo:b1:arith:1}
$1001 = 1 \times 777 + 224$; $777 = 3 \times 224 + 105$; $224 = 2
\times 105 + 14$; $105 = 7 \times 14 + 7$; $14 = 2 \times 7 + 0$. Logo,
$\gcd(1001, 777) = 7$. De trás para diante:
\[
7 = 105 - 7 \times 14
= 105 - 7(224 - 2\times 105) = 15 \times 105 - 7 \times 224
\]
\[
= 15(777 - 3\times 224) - 7\times 224 = 15 \times 777 - 52 \times 224
= 15 \times 777 - 52(1001 - 777) = 67 \times 777 - 52 \times 1001 .
\]
Verificação: $67 \times 777 = 52\,059$ e $52 \times 1001 = 52\,052$;
a diferença é $7$. Par de Bézout: $(u, v) = (-52, 67)$ para $1001u + 777v =
7$.
\end{solution}

\begin{solution}{exo:b1:arith:2}
Como $10 \equiv 1 \pmod 9$: $10^k \equiv 1$, de modo que $\sum_k d_k 10^k
\equiv \sum_k d_k \pmod 9$. Como $10 \equiv -1 \pmod{11}$: $10^k
\equiv (-1)^k$, de modo que o inteiro é congruente à soma alternada
$\sum_k (-1)^k d_k$ módulo $11$ (começando pelo algarismo das \emph{unidades}
com sinal $+$).

$123\,456\,789$: soma dos algarismos $45 \equiv 0 \pmod 9$. Soma alternada a partir
das unidades: $9 - 8 + 7 - 6 + 5 - 4 + 3 - 2 + 1 = 5$, de modo que o número é
$\equiv 5 \pmod{11}$.
\end{solution}

\begin{solution}{exo:b1:arith:3}
Euclides: $237 = 2 \times 91 + 55$; $91 = 1 \times 55 + 36$; $55 = 1
\times 36 + 19$; $36 = 1 \times 19 + 17$; $19 = 1 \times 17 + 2$; $17
= 8 \times 2 + 1$. De trás para diante:
\[
1 = 17 - 8\times 2 = 17 - 8(19 - 17) = 9\times 17 - 8\times 19
= 9(36 - 19) - 8\times 19 = 9\times 36 - 17\times 19
\]
\[
= 9\times 36 - 17(55 - 36) = 26\times 36 - 17\times 55
= 26(91 - 55) - 17\times 55 = 26\times 91 - 43\times 55
\]
\[
= 26\times 91 - 43(237 - 2\times 91) = 112 \times 91 - 43 \times 237.
\]
Logo, $91 \times 112 \equiv 1 \pmod{237}$: as soluções são $x \equiv
112 \pmod{237}$. (Verificação: $91 \times 112 = 10\,192 = 43 \times 237 +
1$.)
\end{solution}

\begin{solution}{exo:b1:arith:4}
$\gcd(17, 39) = 1$: Euclides dá $39 = 2\times 17 + 5$, $17 = 3\times
5 + 2$, $5 = 2\times 2 + 1$ e, de trás para diante,
\[
1 = 5 - 2\times 2 = 5 - 2(17 - 3\times 5) = 7\times 5 - 2\times 17
= 7(39 - 2\times 17) - 2\times 17 = 7\times 39 - 16\times 17 .
\]
Solução particular $(x_0, y_0) = (-16, 7)$. Solução geral da
equação homogênea $17x + 39y = 0$: $x = 39k$, $y = -17k$ (pois
$17 \mid 39y$ e $\gcd(17,39) = 1$ forçam $17 \mid y$ ---
lema de Gauss). Portanto,
\[
(x, y) = (-16 + 39k,\; 7 - 17k), \qquad k \in \Z .
\]
Para o lado direito $5$, multiplique a solução particular por $5$:
$(x, y) = (-80 + 39k,\; 35 - 17k)$, $k \in \Z$.
\end{solution}

\begin{solution}{exo:b1:arith:5}
Para todo \omterm{def:b1:arith:prime}{primo} $p$, com $\alpha = v_p(a)$ e $\beta = v_p(b)$:
\[
v_p\bigl(\gcd(a,b)\bigr) + v_p\bigl(\operatorname{lcm}(a,b)\bigr)
= \min(\alpha, \beta) + \max(\alpha, \beta)
= \alpha + \beta = v_p(ab) .
\]
Dois inteiros positivos com a mesma valoração em todo \omterm{def:b1:arith:prime}{primo} são iguais
(\cref{prop:b1:arith:valuation}), de modo que $\gcd(a,b)\operatorname{lcm}(a,b)
= ab$.
\end{solution}

\begin{solution}{exo:b1:arith:6}
Valorações: $\gcd(a, b) = 2^{\min(10,6)} 3^{\min(4,7)} 5^{\min(2,0)}
7^{\min(0,1)} = 2^6\, 3^4 = 5184$;
$\operatorname{lcm}(a,b) = 2^{10}\, 3^7\, 5^2\, 7$.

Contagem de divisores: um divisor positivo de $n = \prod p_i^{\alpha_i}$ é
exatamente uma escolha $\prod p_i^{\beta_i}$ com $0 \leq \beta_i \leq
\alpha_i$ (\cref{prop:b1:arith:valuation}); as escolhas são
independentes, de modo que há $\prod_i (\alpha_i + 1)$ divisores. Para $a$:
$(10+1)(4+1)(2+1) = 165$.
\end{solution}

\begin{solution}{exo:b1:arith:7}
Suponha $\sqrt p = \frac rq$ com $r, q \in \N^*$, isto é, $p q^2 =
r^2$. Aplique $v_p$: $v_p(pq^2) = 1 + 2v_p(q)$ é ímpar, ao passo que $v_p(r^2)
= 2 v_p(r)$ é par. Um inteiro não pode ter ao mesmo tempo valoração
$p$-ádica ímpar e par: contradição. Logo, $\sqrt p \notin \Q$.
\end{solution}

\begin{solution}{exo:b1:arith:8}
\emph{Fato geral.} Com $\gcd(m,n) = 1$, Bézout dá $mu + nv = 1$.
Ponha $x_0 = b\,mu + a\,nv$. Então $x_0 \equiv a\,nv \equiv a(1 - mu)
\equiv a \pmod m$ e, do mesmo modo, $x_0 \equiv b \pmod n$: existência. Se
$x$ e $x'$ são duas soluções, $m$ e $n$ \omterm{def:b1:arith:divides}{dividem} $x - x'$, de modo que
$mn \mid x - x'$ (\cref{thm:b1:arith:gauss}~(2)): unicidade módulo
$mn$.

\emph{Numericamente:} $m = 7$, $n = 11$: $7 \times (-3) + 11 \times 2 =
1$. Logo, $x_0 = 5 \times 7 \times (-3) + 2 \times 11 \times 2 = -105 +
44 = -61 \equiv 16 \pmod{77}$. Verificação: $16 = 2\times 7 + 2 \equiv 2
\pmod 7$; $16 = 11 + 5 \equiv 5 \pmod{11}$. Soluções: $x \equiv 16
\pmod{77}$.
\end{solution}

\begin{solution}{exo:b1:arith:9}
\omterm{def:b1:complex:field}{Módulo} $7$: Fermat dá $3^6 \equiv 1$, e $1000 = 6 \times 166 + 4$,
de modo que $3^{1000} \equiv 3^4 = 81 \equiv 4 \pmod 7$.

Últimos dois algarismos de $7^{100}$: trabalhe módulo $4$ e módulo $25$.
\omterm{def:b1:complex:field}{Módulo} $4$: $7
\equiv -1$, de modo que $7^{100} \equiv 1$. \omterm{def:b1:complex:field}{Módulo} $25$: $7^2 = 49 \equiv -1$, de modo que
$7^4 \equiv 1$ e $7^{100} = (7^4)^{25} \equiv 1$. Pelo teorema chinês
do resto (\cref{exo:b1:arith:8}), $7^{100} \equiv 1
\pmod{100}$: os dois últimos algarismos são $01$.
\end{solution}

\begin{solution}{exo:b1:arith:10}
Escreva $m = nq + r$, $0 \leq r < n$. Então
\[
2^m - 1 = 2^r\bigl(2^{nq} - 1\bigr) + 2^r - 1,
\]
e $2^n - 1$ \omterm{def:b1:arith:divides}{divide} $2^{nq} - 1 = (2^n - 1)(2^{n(q-1)} + \dots +
1)$. Assim, módulo $2^n - 1$, $\;2^m - 1 \equiv 2^r - 1$ e, como $0 \leq
2^r - 1 < 2^n - 1$, este \emph{é} o resto euclidiano.

Portanto, o \omterm{met:b1:arith:euclid}{algoritmo de Euclides} sobre o par $(2^m - 1, 2^n - 1)$
espelha, expoente por expoente, o algoritmo sobre $(m, n)$: cada
passo de divisão substitui $(m, n)$ por $(n, r)$ em cima e $(2^m - 1,
2^n - 1)$ por $(2^n - 1, 2^r - 1)$ embaixo. O algoritmo de cima
termina em $\gcd(m,n)$, de modo que o de baixo termina em
$2^{\gcd(m,n)} - 1$.
\end{solution}

\begin{solution}{exo:b1:arith:11}
Resolva primeiro $x^2 \equiv 1 \pmod p$: $p \mid (x-1)(x+1)$, de modo que, pelo
lema de Euclides, $x \equiv 1$ ou $x \equiv -1 \pmod p$.

No produto $(p-1)! = 1 \times 2 \times \dots \times (p-1)$, todo
fator $a$ é invertível módulo $p$, e a sua inversa $a^{-1}$ é de novo
um dos fatores (\cref{prop:b1:arith:invmod}). Emparelhe cada $a$ com
$a^{-1}$: os pares multiplicam-se dando $1$, salvo os fatores
emparelhados consigo mesmos ($a = a^{-1}$, isto é, $a^2 \equiv 1$), que ficam
sozinhos --- e estes são exatamente $1$ e $p - 1$. Portanto,
\[
(p-1)! \equiv 1 \times (p - 1) \equiv -1 \pmod p .
\]
(Para $p = 2$: $1! = 1 \equiv -1 \pmod 2$; o argumento de emparelhamento
degenera, mas o resultado vale.)

\emph{Recíproca.} Seja $n \geq 2$ composto, $n = ab$ com $1 < a
\leq b < n$. Se $a < b$, ambos aparecem como fatores distintos de $(n-1)!$,
de modo que $n \mid (n-1)!$ e $(n-1)! \equiv 0 \not\equiv -1$. Se $a = b$
(isto é, $n = a^2$): para $a \geq 3$, tanto $a$ quanto $2a$ são $< n$, logo
$n = a^2 \mid a \times 2a \mid (n-1)!$, mesma conclusão; para $n = 4$,
$(n-1)! = 6 \equiv 2 \not\equiv -1 \pmod 4$.
\end{solution}

\begin{solution}{exo:b1:arith:12}
\begin{enumerate}
  \item Indução. Para $n = 1$: $F_0 = 3 = F_1 - 2 = 5 - 2$.
        Supondo $F_0\cdots F_{n-1} = F_n - 2$:
        \[
        F_0 \cdots F_n = (F_n - 2)F_n
        = \bigl(2^{2^n} - 1\bigr)\bigl(2^{2^n} + 1\bigr)
        = 2^{2^{n+1}} - 1 = F_{n+1} - 2 .
        \]
  \item Sejam $m < n$ e $d = \gcd(F_m, F_n)$. Por (1), $F_m$
        \omterm{def:b1:arith:divides}{divide} $F_n - 2$, de modo que $d$ \omterm{def:b1:arith:divides}{divide} tanto $F_n$ quanto $F_n -
        2$ e, portanto, \omterm{def:b1:arith:divides}{divide} $2$. Mas todo número de Fermat é ímpar,
        logo $d = 1$.
  \item Cada $F_n \geq 3$ tem um divisor \omterm{def:b1:arith:prime}{primo} $p_n$
        (primeiro passo do \cref{thm:b1:arith:euclidprimes}). Se $m
        \neq n$, então $p_m \neq p_n$, pois um \omterm{def:b1:arith:prime}{primo} comum
        dividiria $\gcd(F_m, F_n) = 1$. A \omterm{def:b1:logic:map}{aplicação} $n \mapsto p_n$ é,
        portanto, \omterm{def:b1:logic:inj}{injetiva} de $\N$ nos \omterm{def:b1:arith:prime}{primos}: existem
        infinitos \omterm{def:b1:arith:prime}{primos}.
\end{enumerate}
\end{solution}

\begin{solution}{pb:b1:arith:1}
\textbf{1.} $10! = 3\,628\,800$: dois zeros finais. Valorações
fator a fator: as potências de $2$ vêm de $2, 4 = 2^2, 6, 8 = 2^3,
10$, totalizando $v_2(10!) = 1 + 2 + 1 + 3 + 1 = 8$; as potências de $5$
vêm de $5$ e $10$: $v_5(10!) = 2$. Zeros finais $=
\min(v_2, v_5) = 2$, coerente.

\textbf{2.} Escreva a divisão euclidiana $\lfloor x\rfloor = nq +
r$, $0 \leq r \leq n - 1$. Então $x = nq + r + \{x\}$ com $0 \leq
r + \{x\} < n$, de modo que $\lfloor x/n \rfloor = q = \bigl\lfloor \lfloor
x \rfloor / n \bigr\rfloor$.

\textbf{3.} Seja $\alpha = v_p(a) \leq \beta = v_p(b)$ (troque se
necessário) e escreva $a = p^\alpha a'$, $b = p^\beta b'$ com $p
\nmid a', b'$. Então $a + b = p^\alpha\bigl(a' + p^{\beta -
\alpha}b'\bigr)$, de modo que $v_p(a + b) \geq \alpha = \min$. Se $\alpha <
\beta$, o parêntese vale $a' + p^{\beta-\alpha}b' \equiv a'
\not\equiv 0 \pmod p$: a valoração é exatamente $\alpha$.

\textbf{4.} Os múltiplos de $m$ em $\intint1n$ são $m, 2m, \dots,
qm$, em que $q$ é o maior inteiro com $qm \leq n$, isto é, $q =
\lfloor n/m \rfloor$.

\textbf{5.} Pela fatoração única, $v_p(n!) = \sum_{j=1}^{n}
v_p(j)$. Conte de outro modo: cada $j$ contribui com $v_p(j) =
\#\{k \geq 1 : p^k \mid j\}$, de modo que
\[
v_p(n!) = \sum_{j=1}^n \#\{k : p^k \mid j\}
= \sum_{k\geq1} \#\{j \leq n : p^k \mid j\}
= \sum_{k\geq1} \Bigl\lfloor \frac n{p^k} \Bigr\rfloor
\]
pela questão 4 --- a fórmula de Legendre. A soma é finita: os termos
com $p^k > n$ se anulam.

\textbf{6.} $v_5(1000!) = 200 + 40 + 8 + 1 = 249$ (divisões por
$5, 25, 125, 625$); $v_2(1000!) = 500 + 250 + 125 + 62 + 31 + 15 +
7 + 3 + 1 = 994$. Zeros finais de $1000!$: cada zero consome
um $2$ e um $5$, de modo que há $\min(994, 249) = 249$ deles.

\textbf{7.} Com $n = \sum_i a_ip^i$, a questão 2 dá $\lfloor
n/p^k \rfloor = \sum_{i \geq k} a_ip^{i-k}$ (trunque a expansão na base
$p$). Somando sobre $k \geq 1$ e trocando as duas somas finitas:
\[
v_p(n!) = \sum_{i\geq1} a_i \sum_{k=1}^{i} p^{i-k}
= \sum_{i\geq0} a_i\,\frac{p^i - 1}{p - 1}
= \frac{n - s_p(n)}{p - 1} .
\]

\textbf{8.} Para $p = 2$: $v_2(n!) = n - s_2(n)$. Como $n \geq 1$
tem $s_2(n) \geq 1$, sempre $v_2(n!) \leq n - 1 < n$: $2^n \nmid
n!$. E $v_2(n!) = n - 1$ se, e somente se, $s_2(n) = 1$, se, e somente se, $n$ é potência de
$2$.

\textbf{9.} $n$ tem $\lfloor \log_p n \rfloor + 1$ algarismos na base
$p$, cada um no máximo $p - 1$, de modo que $1 \leq s_p(n) \leq
(p-1)\bigl(\log_p(n) + 1\bigr)$. Substituindo na questão 7:
\[
\frac n{p-1} - \log_p(n) - 1 \;\leq\; v_p(n!) \;<\; \frac n{p-1},
\]
e, dividindo por $n$: $\frac{v_p(n!)}n \to \frac1{p-1}$.

\textbf{10.} $Z(n) - Z(n-1) = v_5(n!/(n-1)!) = v_5(n)$: a contagem
de zeros finais salta de $v_5(n)$ em cada múltiplo de $5$ e é
constante entre eles. $Z(24) = \lfloor24/5\rfloor = 4$ e $Z(25) =
5 + 1 = 6$: em $n = 25$ a contagem salta de $4$ direto para $6$
($v_5(25) = 2$) e, como $Z$ é não decrescente com $Z \leq 4$
antes e $Z \geq 6$ depois, o valor $5$ nunca é atingido: nenhum
fatorial termina em exatamente cinco zeros.

\textbf{11.} Escreva $x = \lfloor x\rfloor + \{x\}$: $\lfloor x +
y\rfloor = \lfloor x\rfloor + \lfloor y\rfloor + \lfloor \{x\} +
\{y\}\rfloor$, e $0 \leq \{x\} + \{y\} < 2$ faz com que o último piso valha
$0$ ou $1$. Então, aplicando Legendre três vezes,
\[
v_p\binom{m+n}m = v_p\bigl((m{+}n)!\bigr) - v_p(m!) - v_p(n!)
= \sum_{k\geq1}\Bigl(
\Bigl\lfloor\frac{m+n}{p^k}\Bigr\rfloor
- \Bigl\lfloor\frac{m}{p^k}\Bigr\rfloor
- \Bigl\lfloor\frac{n}{p^k}\Bigr\rfloor\Bigr),
\]
uma soma finita de $0$s e $1$s (aplique a primeira afirmação a $x =
m/p^k$, $y = n/p^k$).

\textbf{12.} Fixe $k \geq 1$ e escreva $m = p^km_1 + m_0$, $n =
p^kn_1 + n_0$ com $0 \leq m_0, n_0 < p^k$ (divisão euclidiana:
$m_0$ é o número formado pelos $k$ algarismos baixos de $m$). Então
\[
\Bigl\lfloor\frac{m+n}{p^k}\Bigr\rfloor
- \Bigl\lfloor\frac m{p^k}\Bigr\rfloor
- \Bigl\lfloor\frac n{p^k}\Bigr\rfloor
= \Bigl\lfloor\frac{m_0 + n_0}{p^k}\Bigr\rfloor ,
\]
que vale $1$ se $m_0 + n_0 \geq p^k$ e $0$ caso contrário. Mas $m_0 +
n_0 \geq p^k$ diz precisamente que somar os $k$ algarismos baixos de $m$
e de $n$ transborda para a posição $k$ --- um transporte para a posição
$k$ no algoritmo escolar da adição. Somando sobre $k$:
$v_p\binom{m+n}m$ é o número de transportes na
adição $m + n$ na base $p$. (Kummer, 1852.)

\textbf{13.} Aplique Kummer a $m = j$, $n = p^k - j$, cuja soma é $p^k =
(1\underbrace{0\cdots0}_{k})_p$. Seja $a = v_p(j)$, de modo que os
algarismos de $j$ na base $p$ nas posições $0, \dots, a-1$ são $0$ e o
algarismo na posição $a$ é não nulo. Os algarismos de $p^k - j$ abaixo da
posição $a$ também são $0$ ($p^k - j = p^a(p^{k-a} - j/p^a)$). Na posição
$a$, os dois algarismos não nulos devem somar $p$ (algarismo resultante
$0$): um transporte; em cada posição $a+1, \dots, k-1$, os algarismos
mais o transporte que chega somam $p$ (de novo, algarismo resultante
$0$): o transporte se propaga. Total: $k - a$ transportes, de modo que
$v_p\binom{p^k}j = k -
v_p(j)$. Para $k = 1$: $v_p\binom pj = 1$ para $0 < j < p$, a
\omterm{def:b1:arith:divides}{divisibilidade} usada no \cref{thm:b1:arith:fermat}.

\textbf{14.} Pela forma digital (questão 7), usando $s_2(2n) =
s_2(n)$ (basta acrescentar um algarismo zero):
\[
v_2\binom{2n}n = \bigl(2n - s_2(2n)\bigr) - 2\bigl(n -
s_2(n)\bigr) = 2s_2(n) - s_2(2n) = s_2(n) \geq 1 :
\]
$\binom{2n}n$ é sempre par, e $v_2 = 1$ (isto é, $\binom{2n}n
\equiv 2 \pmod 4$) exatamente quando $s_2(n) = 1$, ou seja, quando $n$ é
uma potência de $2$.

\textbf{15.} Vandermonde com $m = n = k = p$: $\binom{2p}p =
\sum_{j=0}^p \binom pj\binom p{p-j} = \sum_{j=0}^p \binom pj^2$.
Para $0 < j < p$, $p \mid \binom pj$ (questão 13), de modo que $\binom
pj^2 \equiv 0 \pmod p$; os termos das pontas dão $1 + 1$:
$\binom{2p}p \equiv 2 \pmod p$.

\textbf{16.} Base $3$: $500 = 486 + 9 + 3 + 2$, algarismos (do menos para o
mais significativo) $(2, 1, 1, 0, 0, 2)$, de modo que $s_3(500) = 6$; e $1000 = 729 +
243 + 27 + 1$, algarismos $(1, 0, 0, 1, 0, 1, 1)$, de modo que $s_3(1000) = 4$.
\emph{Kummer:} some $500 + 500$ na base $3$: posição $0$: $2 + 2 =
4$, algarismo $1$, transporte $1$; posição $1$: $1 + 1 + 1 = 3$, algarismo $0$,
transporte $1$; posição $2$: $1 + 1 + 1 = 3$, algarismo $0$, transporte $1$;
posição $3$: $0 + 0 + 1 = 1$, sem transporte; posição $4$: $0$;
posição $5$: $2 + 2 = 4$, algarismo $1$, transporte $1$; posição $6$:
o transporte cai ali: algarismo $1$. Quatro transportes: $v_3\binom{1000}{500} = 4$.
\emph{Legendre:} $v_3(1000!) = \frac{1000 - 4}2 = 498$ e
$v_3(500!) = \frac{500 - 6}2 = 247$, de modo que $v_3\binom{1000}{500} =
498 - 2\times247 = 4$. Os dois cálculos concordam --- e os
algarismos da adição $(1, 0, 0, 1, 0, 1, 1)$ reproduzem $1000$, como
devem.

\textbf{17.} Por Kummer ($p = 2$, $m = k$, $n' = n - k$):
$\binom nk$ é ímpar se, e somente se, a adição $k + (n - k)$ na base $2$ não
tem transporte, isto é, se, e somente se, em cada posição os algarismos
satisfazem $k_i + (n -
k)_i = n_i$; nesse caso, $k_i \leq n_i$ para todo $i$. Reciprocamente,
se $k_i \leq n_i$ para todo $i$, então o número de algarismos $n_i -
k_i$ é $n - k$ e a adição não tem transportes. Mesma demonstração na
base $p$: $p \nmid \binom nk$ se, e somente se, todo algarismo de $k$ na
base $p$ é no máximo o algarismo correspondente de $n$.

\textbf{18.} Contando os $k \in \intint0n$ cujos algarismos obedecem a
$k_i \leq n_i$: cada algarismo de $k$ é escolhido independentemente entre
$n_i + 1$ valores, o que dá $\prod_i (n_i + 1)$ escolhas; na base $2$
isso é $2^{\#\{i : n_i = 1\}} = 2^{s_2(n)}$. Linha $4 =
(100)_2$: $2^1 = 2$ entradas ímpares --- de fato, $1, 4, 6, 4, 1$ tem entradas
ímpares apenas nas pontas. Linha $5 = (101)_2$: $2^2 = 4$ --- de fato,
$1, 5, 10, 10, 5, 1$.

\textbf{19.} Todas as entradas interiores são pares $\iff$ a linha tem
exatamente $2$ entradas ímpares (as duas pontas são sempre ímpares)
$\iff 2^{s_2(n)} =
2 \iff s_2(n) = 1 \iff n$ é uma potência de $2$.

\textbf{20.} Na soma da questão 11, o $k$-ésimo termo se anula assim
que $p^k > m + n$ (os três pisos são então iguais; com efeito, o
primeiro vale $0$ quando $p^k > m+n$; mais simplesmente, cada termo é $0$).
Portanto, no máximo $\lfloor \log_p(m+n)\rfloor$ termos são não nulos, cada
um valendo $1$: $a = v_p\binom{m+n}m \leq \log_p(m+n)$, isto é, $p^a \leq
m + n$.

\textbf{21.} Para todo \omterm{def:b1:arith:prime}{primo} $p$, $v_p\bigl(\operatorname{lcm}(1,
\dots, 2n)\bigr) = \lfloor\log_p(2n)\rfloor$ (a maior potência de
$p$ que não excede $2n$ aparece entre $1, \dots, 2n$). A questão 20
com $m = n$ dá $v_p\binom{2n}n \leq \lfloor\log_p(2n)\rfloor$
para todo $p$: pela \cref{prop:b1:arith:valuation}, $\binom{2n}n
\mid \operatorname{lcm}(1, \dots, 2n)$. Quanto ao tamanho: a razão
$\binom{2n}{k+1}/\binom{2n}k = \frac{2n-k}{k+1} \geq 1$ exatamente
para $k < n$, de modo que a entrada central é a maior das $2n + 1$
entradas da linha $2n$, donde $4^n = \sum_k \binom{2n}k \leq
(2n+1)\binom{2n}n$. Combinando:
\[
\operatorname{lcm}(1, \dots, 2n) \geq \binom{2n}n \geq
\frac{4^n}{2n + 1} .
\]
Se houvesse poucos \omterm{def:b1:arith:prime}{primos} abaixo de $2n$, o mmc não poderia ser tão
grande: o crescimento exponencial do mmc é um traço quantitativo da
abundância dos \omterm{def:b1:arith:prime}{primos}.

\textbf{22.} $Z(n) = \sum_k\lfloor n/5^k\rfloor \approx \frac
n4$, de modo que se deve mirar perto de $n = 4 \times 2026 = 8104$: $Z(8104) = 1620 +
324 + 64 + 12 + 2 = 2022$. Suba de $5$ em $5$: $Z(8110)
= 2024$, $Z(8115) = 2025$, e
\[
Z(8120) = 1624 + 324 + 64 + 12 + 2 = 2026 .
\]
Como $Z$ é constante entre múltiplos de $5$ e $Z(8119) =
Z(8115) = 2025$, o menor $n$ com pelo menos $2026$ zeros
finais é $n = 8120$.

\textbf{23.} Base $7$: $50 = 49 + 1$, algarismos (do menos para o mais
significativo) $(1, 0,
1)$. Somando $50 + 50$: posição $0$: $1 + 1 = 2 < 7$, sem transporte;
posição $1$: $0 + 0 = 0$; posição $2$: $1 + 1 = 2 < 7$, sem
transporte. Sem transportes, de modo que, por Kummer, $v_7\binom{100}{50} = 0$: $7 \nmid
\binom{100}{50}$. Legendre concorda: $v_7(100!) = \lfloor 100/7
\rfloor + \lfloor 100/49 \rfloor = 14 + 2 = 16$ e $v_7(50!) = 7
+ 1 = 8$, de modo que $v_7\binom{100}{50} = 16 - 2\times8 = 0$.

\textbf{24.} (i) A fatoração única sustenta a própria
definição de $v_p$ e a sua aditividade, logo a fórmula de Legendre
e toda conclusão de \omterm{def:b1:arith:divides}{divisibilidade}
(\cref{prop:b1:arith:valuation}). (ii) A divisão euclidiana produziu
a identidade de truncamento da questão 2 e a separação $m = p^km_1 +
m_0$ que isola o transporte (questão 12). (iii) Contagem: a
contagem dos múltiplos de $m$ (questão 4), o produto de escolhas de
algarismos (questão 18) e a estimativa da soma de linha $4^n \leq
(2n+1)\binom{2n}n$ (questão 21) são todos
argumentos ao estilo do \cref{ch:b1:counting}.

\textbf{25.} Legendre converte ``que potência de $p$ \omterm{def:b1:arith:divides}{divide}
$n!$'' em aritmética de algarismos na base $p$; Kummer comprime a
resposta para os \omterm{def:b1:counting:objects}{coeficientes binomiais} nos transportes de uma única
adição --- a \omterm{def:b1:arith:divides}{divisibilidade}, aparentemente uma propriedade global de
números enormes, é lida localmente, algarismo a algarismo. A questão 16
é o paradigma: quatro transportes, calculados à mão, determinam a
potência exata de $3$ num número com centenas de algarismos. E a
questão 21 mostra o mesmo círculo de ideias roçando águas profundas: uma
estimativa inferior exponencial para $\operatorname{lcm}(1, \dots, 2n)$ é
um primeiro passo, inteiramente elementar, rumo ao teorema dos \omterm{def:b1:arith:prime}{números
primos}, cuja demonstração está muito além deste volume. Todo o
instrumental --- divisão, mdc, valorações --- é reencenado para os
polinômios no \cref{ch:b1:poly}, em que o análogo de uma expansão em
algarismos é a expansão em potências de $(X - a)$.
\end{solution}
